\documentclass[11pt,letterpaper]{article}

% Compile with pdfLaTeX (or select pdfLaTeX in Overleaf).
\usepackage[margin=1in]{geometry}
\usepackage[T1]{fontenc}
\usepackage{lmodern}
\usepackage{amsmath,amssymb}
\usepackage{booktabs}
\usepackage{enumitem}
\usepackage[hidelinks]{hyperref}
\usepackage{microtype}

\setlength{\parindent}{0pt}
\setlength{\parskip}{0.5em}
\setlist[enumerate]{leftmargin=*,itemsep=1em}

\newcommand{\studentname}{Jiahao Cheng}
\newcommand{\studentid}{12549801}
\newcommand*{\repourl}{} % Optional URL for the homework_3 directory.
\newcommand{\partlabel}[1]{\par\medskip\textbf{(#1)}\par}
\newcommand{\E}{\mathbb{E}}

\begin{document}

\begin{center}
{\Large\bfseries Machine Learning 1 --- Homework 3}\\[5pt]
Fall 2026 \quad $\vert$ \quad DATA 37711\\[8pt]
\studentname\quad $\vert$ \quad \studentid
\ifx\repourl\empty\else
  \quad $\vert$ \quad \href{\repourl}{Code Repo: homework\_3}
\fi
\end{center}
\hrule
\vspace{8pt}

\section*{Problem 1}

\partlabel{a}
\begin{enumerate}[label=(\roman*),itemsep=0.3em]
    \item Poisson is a reasonable starting point for a count response because its
    support is the set of nonnegative integers. Its rate represents the
    conditional expected count. It also assumes that the conditional variance
    equals the conditional mean.
    \item The ensemble score \(f(x)\) can take any real value. The exponential
    inverse link gives a strictly positive rate \(\lambda(x)=\exp(f(x))\).
    Equivalently, the canonical link is \(f(x)=\log\lambda(x)\).
\end{enumerate}

\partlabel{b}
The Poisson probability mass is
\[
P(Y=y\mid X=x)=\frac{\exp(-\lambda(x))\lambda(x)^y}{y!}.
\]
With \(\lambda(x)=\exp(f(x))\), its negative log-likelihood is
\[
L(y,f(x))=\exp(f(x))-yf(x)+\log(y!).
\]
The final term is constant in \(f\), so it can be dropped when minimizing the
empirical loss.

\partlabel{c}
\[
\frac{\partial L}{\partial f(x)} = \exp(f(x))-y.
\]

\partlabel{d}
Start with a constant score \(f_0\). At step \(j\), compute the negative
gradient \(r_i=y_i-\exp(f_{j-1}(x_i))\). Fit a regression tree to these targets,
\[
h_j \in \operatorname*{arg\,min}_{h\in\mathcal H}
\sum_{i=1}^{n}\bigl(r_i-h(x_i)\bigr)^2.
\]
Choose a step size \(\rho_j\), for example by minimizing the empirical Poisson
loss along \(h_j\), and update \(f_j=f_{j-1}+\rho_j h_j\).

\clearpage
\section*{Problem 2}

\partlabel{a}
Suppose \(k\geq 1\) candidate models are fit using \(n\) training examples.
Conditional on the training data and training randomness, assume the
candidates \(f_1,\ldots,f_k\) are fixed. Let \(Z_1,\ldots,Z_\ell\) be i.i.d.\ test
observations drawn independently of the data and randomness used to construct
the candidates. Assume the per-observation loss lies in \([0,B]\), with
\(B>0\). For candidate \(j\), define
\[
R_j=\E[L(f_j,Z)],\qquad
\widehat R_j=\frac{1}{\ell}\sum_{i=1}^{\ell}L(f_j,Z_i).
\]
Let \(\widehat j\) minimize \(\widehat R_j\) over the \(k\) candidates. For
\(0<\delta<1\), with probability at least \(1-\delta\),
\[
\left|\widehat R_{\widehat j}-R_{\widehat j}\right|
\leq B\sqrt{\frac{\log(2k/\delta)}{2\ell}}.
\]
For any fixed \(j\), Hoeffding gives
\(\Pr(|\widehat R_j-R_j|>\epsilon)
 \leq 2\exp(-2\ell\epsilon^2/B^2)\).
A union bound over the \(k\) candidates yields
\[
\Pr\!\left(\max_{1\leq j\leq k}|\widehat R_j-R_j|>\epsilon\right)
\leq 2k\exp(-2\ell\epsilon^2/B^2).
\]
Set the right-hand side equal to \(\delta\). The resulting event controls all
candidates simultaneously, including the selected candidate. The conditional
probability bound holds for every training realization, so it also holds
unconditionally.

The deviation bound decreases with \(\ell\) and depends logarithmically on
\(k\). The training size \(n\) can affect the candidates' risks, but it does not
enter this conditional deviation bound. For the wine regressors below, both
models predict averages of training responses in \([1,10]\), so the
per-observation squared loss is bounded by \(81\).

\partlabel{b}
The same test set can give an optimistic estimate if it is used to select the
best candidate and to report that candidate's risk. Conditional on training,
\[
\E[\widehat R_{\widehat j}]
=\E[\min_j\widehat R_j]
\leq\min_j\E[\widehat R_j]
=\min_jR_j
\leq\E[R_{\widehat j}].
\]
With many candidates and a small test set, a candidate can win because of
favorable test noise. An independent test set gives an unbiased risk estimate
conditional on the model already selected using validation data.

\clearpage
\section*{Problem 2 (continued)}

\partlabel{c}
I used the public white-wine CSV from the \href{https://doi.org/10.24432/C56S3T}{UCI Wine Quality dataset}
(Cortez et al., 2009), which the assignment links to. The course-provided CSV
was unavailable for comparison. The response was the integer \texttt{quality}
and the other 11 columns were predictors. I used seed 42 for both random splits.
I first held out 20\% for testing, then held out 12.5\% of the remaining rows
for validation. This produced 3,428 training, 490 validation, and 980 test rows.

\textbf{Random forest.}
I varied the number of trees (10, 20, 30, 40, 50, 60), maximum depth
(3, 5, 7, 9, 12, 16, unrestricted), number of candidate predictors per split
(half, square root, or all), and minimum leaf size (1, 2, 4, 8).
These control ensemble size, tree complexity, feature randomness at each
split, and the smallest allowed leaf. I used the original feature values.
For each combination of the other settings, I extended the forest at each
successive tree count while keeping the seed fixed.

\textbf{K-nearest neighbors.}
I varied the number of neighbors (1 through 42), weighting (uniform or
inverse distance), distance metric (Manhattan or Euclidean), and scaling
(standard, robust, or min--max). The count and weights determine how
neighboring responses are averaged. The metric and scaling determine which
points are nearby. I fit each scaler using only the training rows.

Each grid had 504 settings, exceeding the 490 validation rows. I fit each
candidate on the training split, selected the lowest validation MSE within
each model class, and evaluated the winning setting on the held-out test split.

\begin{center}
\begin{tabular}{lrrr}
\toprule
Model & Grid size & Validation MSE & Test MSE \\
\midrule
Random forest & 504 & 0.3081 & 0.3553 \\
K-nearest neighbors & 504 & 0.2987 & 0.3686 \\
\bottomrule
\end{tabular}
\end{center}

The winning forest used 60 trees, unrestricted depth, the square root of the
feature count per split, and minimum leaf size 1. The winning KNN used
25 neighbors, inverse-distance weighting, Manhattan distance, and standard scaling.

For this split, test MSE exceeded validation MSE by 0.0472 for the forest and
0.0698 for KNN. KNN ranked first on validation, while the forest ranked first
on test. Minimizing validation MSE across the grid can make the winning
validation score optimistic. The independent test MSE also varies across samples.

Documentation:
\href{https://scikit-learn.org/stable/modules/generated/sklearn.ensemble.RandomForestRegressor.html}{RandomForestRegressor}
and
\href{https://scikit-learn.org/stable/modules/generated/sklearn.neighbors.KNeighborsRegressor.html}{KNeighborsRegressor}.
The UCI data are licensed under \href{https://creativecommons.org/licenses/by/4.0/}{CC BY 4.0}.

\clearpage
\section*{Problem 3}

For each dataset (\texttt{dset0-1.csv} and \texttt{dset1-1.csv}), I used the
first column as the binary response and the other
five columns as predictors. An 80\%/20\% random split with seed 42 gave 400
training and 100 test rows per dataset. I used the same split for all five
values of \texttt{n\_estimators}, with \texttt{GradientBoostingClassifier},
learning rate 0.1, and maximum tree depth 3. Error is the percentage of
incorrect classifications.

\begin{center}
\begin{tabular}{rrrrr}
\toprule
& \multicolumn{2}{c}{\texttt{dset0}} & \multicolumn{2}{c}{\texttt{dset1}} \\
\cmidrule(lr){2-3}\cmidrule(lr){4-5}
Trees & Training & Test & Training & Test \\
\midrule
1     & 42.50\% & 43\% & 43.75\% & 45\% \\
10    & 13.25\% & 27\% & 12.75\% & 23\% \\
100   &  0.50\% & 18\% &  1.25\% & 20\% \\
300   &  0.00\% & 17\% &  0.00\% & 18\% \\
1,000 &  0.00\% & 19\% &  0.00\% & 18\% \\
\bottomrule
\end{tabular}
\end{center}

\partlabel{a}
On \texttt{dset0}, test error fell from 43\% with one tree to a minimum of 17\%
at 300 trees. It rose to 19\% at 1,000 trees while training error stayed at
zero. The increase represents two additional misclassifications on this test
split.

\partlabel{b}
On \texttt{dset1}, test error fell from 45\% with one tree to 18\% at 300 trees.
Both test error and training error stayed at 18\% and 0\%, respectively, at
1,000 trees. The last two settings therefore tied for the lowest observed
test error.

\partlabel{c}
I sorted each dataset by its first predictor and divided the rows into five
equally sized groups. From the lowest to the highest group, the positive-class
rates were 89\%, 63\%, 47\%, 20\%, and 8\% in \texttt{dset1}, compared with
52\%, 37\%, 32\%, 34\%, and 58\% in \texttt{dset0}. A single threshold can
capture part of the monotone pattern in \texttt{dset1}. The higher rates at both
ends in \texttt{dset0} can require multiple splits. This may help explain why
\texttt{dset1} has lower error at 10 trees despite having slightly higher
one-tree error. These marginal patterns alone cannot establish why the
learning curves differ.

\end{document}
